\section{Angels and Demons}

\begin{quotex}
Since we cannot be present at the holy ceremonies in honor of the angels, we should not let this time of devotion go by
fruitlessly; rather, such time as we do not spend in singing their praises, we should spend in writing about them. And
because our aim is to present as best we can the excellence of the holy angels, we ought to begin with
man's earliest conjectures about the angels. In this way, we shall be in a position to accept whatever we
find that agrees with faith, and refute whatever is opposed to Catholic teaching. \flright{\textsc{Thomas Aquinas}, \textit{De
Angelis}}

\end{quotex}
In \textit{The Multiple States of Being}, \textbf{Rene Guenon} recommends Aquinas' short theological treatise on the
angels\footnote{\url{https://www.gornahoor.net/?page_id=13278}} as an introduction to an understanding of angels as higher, supra-individual states of being. Aquinas himself
looks back to the ancient Greeks to begin his study of angelic beings. After all, if angels, demons, and gods are real
entities or states of being and not merely figments of the imagination or objects of belief, then they can be
understood by any metaphysician. Aquinas explains:

\begin{quotex}
Now it is clear that every intellectual substance receives the intellected form according to its totality, or otherwise
it would not be able to know it in its totality. For it is thus that the intellect understands a thing insofar as the
form of that thing exists in it.

\end{quotex}
Thus, a complete understanding of the angelic hierarchy is possible only for the intellect that contains the form, that
is, the idea of the angel, that is, it participates in the higher form. This is the metaphysical principle, “to know is
to be”.

\paragraph{Gods and the Celestial Hierarchy}
Confusion results from terminology, since we tend to consider the “angels” as good and “demons” as bad. However, Aquinas
clarifies this by pointing out that angels can be evil and that according to \textbf{Plato}, demons can be good or
evil. Regarding Plato, Aquinas writes:

\begin{quotex}
In this way, therefore, between us and the highest God, it is clear that they posited four orders, namely, that of the
secondary gods, that of the separate intellects, that of the heavenly souls, and that of the good or wicked demons. If
all these things were true, then all these intermediate orders would be called by us “angels”, for Sacred Scripture
refers to the demons themselves as angels. 

\end{quotex}
So, Aquinas accepted the gods as angels. He also considered that the planets were ensouled. Hence, he seems to be
implying that the gods Mars, Venus, Jupiter, Saturn and so on were actually angels. This aligns him with the Hermetic
Tradition. \textbf{Valentin Tomberg}, in Meditations on the Tarot, assigns the correspondences listed in Table~\ref{tab:AngDem2011}.

\begin{table}[h]
\begin{center}
\begin{tabular}{*2c}\toprule
Archangel &
Planet\\\midrule
Michael &
Sun\\
Gabriel &
Moon\\
Raphael &
Mercury\\
Anael &
Venus\\
Zachariel &
Jupiter\\
Oriphiel &
Saturn\\
Samael &
Mars\\\bottomrule
\end{tabular}
\end{center}
\caption{Angel-Planet correspondence according to Tomberg.}\label{tab:AngDem2011}
\end{table}

\paragraph{Preternatural Phenomena}
Aquinas recognizes certain preternatural phenomena and attributes them to the influence of higher intellectual
substance. Keep in mind that in the following passage, the Platonic “demon” is the same as the Catholic “angel”, so no
inference should be made that it necessarily implies an evil or satanic influence.

\begin{quotex}
[Certain followers of Aristotle] say that it is through the influence of the stars that persons who are possessed
sometimes foretell future events, for the realization of which there is a certain disposition in nature through the
heavenly bodies. But in such cases, there are manifestly certain works which cannot in any way be reduced to a
corporeal cause. For example, that people in a trance should speak in a cultivated way of sciences which they do not
know, since they are unlettered folk; and that those who have scarcely left the village in which they were born, speak
with fluency the vernacular of a foreign people. Likewise, in the works of magicians, certain images are said to be
conjured up which answer questions and move about, all of which could not be accomplished by any corporeal cause.
Therefore, as the Platonists see it, who could evidently assign a cause of these effects, except to say that these are
brought about through demons.

\end{quotex}
\paragraph{On Infallibility}
We see here that Aquinas asserts the traditional teaching that gnosis is infallible. Unfortunately, self-deception is
not so easily overcome.

\begin{quotex}
Accordingly, Augustine says in the Book of Eighty-Three Questions: “Everyone who is deceived, that, indeed, in which he
is deceived, he does not understand.” And accordingly, concerning those things which we grasp properly by our intellect
as well as concerning the first principles, no one can be deceived.

\end{quotex}

\flright{\itshape Posted on 2011-08-31 by Cologero}

\begin{center}* * *\end{center}


\begin{footnotesize}\begin{sffamily}

\texttt{Perennial on 2011-08-31 at 03:45 said: }

Based on similar analysis, St. Thomas conceded the possibility of legitimate astrology, practiced within certain limits.
I could see where the information contained here would dovetail very nicely with St. Thomas' analysis on
that question. It would be interesting to see St. Thomas' ideas on the relation between the heavenly bodies
and the spirit explored in more depth.

\hfill

\texttt{Charlotte Cowell on 2011-08-31 at 18:24 said: }

The daemon of Socrates was indeed instructive, while the demons of the modern world are known to serve the purpose of
testing the motives and faith of humans. The Lesser Guardian of the Threshold springs to mind here. For a Christian,
though, I wonder if Hermes, for instance, can truly be considered equivalent to (for example), Michael or Gabriel, as
there seems to be a major difference in function. Let me be blunt to illustrate this point. One can imagine taking (or
being taken by) Hermes as a lover – as nymphs and goddesses alike would attest – but Michael? No, I don't
think so….But there are also definite similarities because both Hermes and Gabriel are heraldic messengers. Then there
is the question of Nephilim that crops up for the occultist. There is a clear distinction between a Nephilim and an
Elohim – should we equate it to the difference between ‘good' and
‘evil' or ‘light' and `dark’, I wonder?
Certainly there is a difference in quality. But perhaps it is analogous to the difference in quality between a rain
cloud and the wispy clouds you see in the sun. The robe of one would feel like petals and the silk of cobwebs – the
finest possible substance one can possibly imagine, the way that light feels. The robe of another would feel heavier,
like storm clouds….one would bring you absolute joy and freedom, restoring you to a state of absolute innocence and
joy, whereas another would hold you in suspense, hypnotically, hardly daring to breathe or blink lest the eye of the
Watcher should fall upon the self. Who could bear to be faced with that terrible gaze? Guardian angels, in the
meantime, are always loving and protective, best not forget them!


\hfill

\texttt{logres on 2011-08-31 at 22:36 said: }

Funny you should mention the Lesser Guardian of the Threshold – I just came across Tomasso Palamadessi.


\hfill

\texttt{Charlotte Cowell on 2011-09-01 at 03:44 said: }

i had never heard of him until you mentioned his name but – having just looked briefly into it – he seems like a very
interesting character, i shall investigate further!

\hfill
\end{sffamily}\end{footnotesize}